\documentclass{amsart}
\usepackage{amsmath,amssymb,amsthm,hyperref}
\usepackage{mathtools}

\newtheorem{theorem}{Theorem}
\newtheorem{lemma}{Lemma}
\newtheorem{proposition}{Proposition}
\newtheorem{corollary}{Corollary}
\theoremstyle{definition}
\newtheorem{definition}{Definition}
\newtheorem{remark}{Remark}

\newcommand{\E}{\mathbb{E}}
\newcommand{\R}{\mathbb{R}}
\newcommand{\kB}{k_{B}}
\newcommand{\Wext}{\langle W_{\mathrm{ext}}\rangle}
\newcommand{\dd}{\mathrm{d}}
\DeclareMathOperator{\KL}{D}

\begin{document}

\title{Stochastic thermodynamics of feedback flashing ratchets:
information bound, maximal entropy reduction, and tight work/power/efficiency bounds}
\author{}
\date{}
\maketitle

\section{Set-up and standing assumptions}

We make the dynamical model precise so that every later statement is
unambiguous.

\paragraph{Particle dynamics.}
A Brownian particle has position $X_t\in\R$ obeying the overdamped Langevin
equation
\begin{equation}\label{eq:langevin}
\gamma\,\dd X_t = -\,a_t\,V'(X_t)\,\dd t + \sqrt{2\gamma \kB T}\;\dd B_t ,
\end{equation}
where $\gamma>0$ is the friction coefficient, $T>0$ the bath temperature,
$B_t$ a standard Wiener process, and
\[
V:\R\to\R,\qquad V(x+L)=V(x),\qquad V\in C^{1},
\]
a fixed $L$-periodic asymmetric potential. The \emph{control variable}
\[
a_t\in\{0,1\}
\]
is piecewise constant in $t$: $a_t=1$ means the potential is ON,
$a_t=0$ means it is OFF (free diffusion). Equation \eqref{eq:langevin} is
understood in the It\^o sense; existence and uniqueness of a strong solution
on each interval where $a_t$ is constant follow from the Lipschitz/linear-growth
property of $V'$ implied by $V\in C^1$ and periodicity (so $V'$ is bounded).

\paragraph{Feedback protocol.}
Let $[0,\tau]$ be the operating window. A \emph{feedback protocol}
$\pi$ is a rule that produces the control trajectory
$a_{[0,\tau]}:=(a_t)_{t\in[0,\tau]}$ as a (possibly randomised) functional
of measurements of the particle. We treat the two standard cases jointly.

\begin{itemize}
\item \emph{Discrete-time feedback.} At sampling instants
$0\le t_1<t_2<\dots<t_N\le\tau$ the controller obtains measurement outcomes
$Y_k=m(X_{t_k},\xi_k)$, where $m$ is a measurement channel and $\xi_k$ is
internal measurement noise (so $Y_k$ may be noisy/coarse-grained). After the
$k$-th measurement the controller chooses the control value on
$[t_k,t_{k+1})$ as a measurable, possibly randomised function of the
\emph{record} $Y^{(k)}:=(Y_1,\dots,Y_k)$. Thus $a_t$ on $[t_k,t_{k+1})$ is a
function of $Y^{(k)}$.

\item \emph{Continuous-time feedback.} The controller has access to the
filtration $\mathcal F^Y_t=\sigma(Y_s:s\le t)$ generated by a (possibly noisy)
observation process $Y_s$, and $a_t$ is $\mathcal F^Y_t$-adapted (causal,
nonanticipating).
\end{itemize}

In both cases the control sequence is \emph{non-Markovian} in the sense that
$a_t$ generally depends on the whole past record, not only on the current
state.

\paragraph{Thermodynamic quantities.}
For a control trajectory $a_{[0,\tau]}$ the stochastic work injected by the
external agent that flashes the potential is, by the first law for the
flashing ratchet (the potential changes only at the switching instants, where
the particle position is frozen),
\begin{equation}\label{eq:work}
W \;=\; \sum_{\text{switch times }s} \bigl(a_{s^+}-a_{s^-}\bigr)\,V(X_s)
\;=\; \int_0^\tau V(X_t)\,\dd a_t ,
\end{equation}
the last expression being a Stieltjes integral against the pure-jump process
$a_t$. The work \emph{extracted} is $W_{\mathrm{ext}}=-W$. The free energy of
the ON state at canonical equilibrium is
$F=-\kB T\log\!\int_0^L e^{-V(x)/\kB T}\dd x$ (per period), and $\Delta F$
denotes the change of free energy between the initial and final ensembles.

Throughout, $\langle\cdot\rangle$ denotes expectation over both the thermal
noise $B$ and the measurement/protocol randomness.

\section{(a) The information quantity for an arbitrary protocol}

\subsection{Reference (open-loop) dynamics and the path measure}
Fix the operating window and an initial law $\rho_0$ for $X_0$. A control
trajectory $a_{[0,\tau]}$ together with the Langevin dynamics
\eqref{eq:langevin} defines a probability measure on particle paths. The
crucial object is the joint law of \emph{(particle path, control record)}.

Let $\Omega$ be the path space of $X$, let $\mathcal Y$ be the space of
measurement records (the vector $Y^{(N)}$ in discrete time, or the path
$Y_{[0,\tau]}$ in continuous time), and let
\[
P(\dd x,\dd y)
\]
be the joint law generated by the actual feedback protocol $\pi$.

We compare $\pi$ to a \emph{reference protocol} that produces an identically
distributed control trajectory but \emph{without using the measurement to
decide it}: i.e.\ the control is drawn from the same marginal law of
$a_{[0,\tau]}$, but \emph{independently} of $X$. Concretely, let
$P_X(\dd x)$ be the marginal law of the particle path and $P_Y(\dd y)$ the
marginal law of the control record under $\pi$. The reference (``no
feedback / blind replay'') joint law is the product
\[
P^{\mathrm{ref}}(\dd x,\dd y) := P_X(\dd x)\,P_Y(\dd y).
\]

\subsection{Definition of $I$}

\begin{definition}[Information of a feedback protocol]\label{def:I}
For a feedback protocol $\pi$ on the flashing ratchet over $[0,\tau]$, the
\emph{information} is the mutual information between the particle path and the
controller's record,
\begin{equation}\label{eq:Idef}
I \;:=\; \mathcal I\!\bigl(X_{[0,\tau]};\,Y\bigr)
\;=\; \KL\!\bigl(P \,\big\|\, P^{\mathrm{ref}}\bigr)
\;=\; \E\!\left[\log\frac{\dd P}{\dd P^{\mathrm{ref}}}(X,Y)\right]\ \ge 0,
\end{equation}
where $\KL(\cdot\|\cdot)$ is relative entropy and the equality with the
Kullback--Leibler divergence to the product of marginals is the defining
identity of mutual information.
\end{definition}

For \emph{discrete-time} protocols, $Y=Y^{(N)}=(Y_1,\dots,Y_N)$ and the chain
rule for mutual information together with causality of the control gives the
operationally transparent decomposition
\begin{equation}\label{eq:Ichain}
I=\sum_{k=1}^{N}\mathcal I\!\bigl(X_{[0,\tau]};\,Y_k \,\big|\, Y^{(k-1)}\bigr).
\end{equation}
This is exactly the \emph{transfer-entropy-type} sum appropriate to a
non-Markovian sequence of control actions: at the $k$-th step the controller
extracts $\mathcal I(X;Y_k\mid Y^{(k-1)})$ nats of \emph{new} information about
the trajectory, conditioned on everything learned so far. For
\emph{continuous-time} protocols, \eqref{eq:Idef} is the continuous-time
limit of \eqref{eq:Ichain}, i.e.\ a transfer-entropy \emph{rate} integrated
over $[0,\tau]$; when the observation is the full position with measurement
noise it is the directed-information / transfer entropy from $X$ to the
control.

\begin{remark}
$I$ depends only on the joint law of (trajectory, record) and is
manifestly nonnegative, finite when the measurement has finitely many
outcomes per step, and reduces to the ordinary Shannon mutual information for
a single measurement. It is the quantity appearing in the generalised second
law: for any feedback protocol on \eqref{eq:langevin},
\begin{equation}\label{eq:2ndlaw}
\Wext \;\le\; -\Delta F + \kB T\, I .
\end{equation}
We do not re-derive \eqref{eq:2ndlaw} (it is the Sagawa--Ueda generalised
second law, valid for any measurement-and-feedback protocol on a Langevin
system); we take it as the framework hypothesis stated in the problem and
\emph{compute the ratchet-specific value of $I$ and its maximum}. The
saturating condition of \eqref{eq:2ndlaw} is reversibility of the controlled
dynamics; we use it in \S4 to identify saturating protocols.
\end{remark}

\section{(b) Maximal entropy reduction $I_{\max}$}

The entropy production of the controlled system is bounded below by $-I$, so
the maximal \emph{entropy reduction} achievable by feedback equals the maximal
value of $I$ over admissible protocols. The decisive structural fact about the
flashing ratchet is that the controller's \emph{action space is binary}:
$a_t\in\{0,1\}$. We exploit this to bound $I$ by the entropy of the control
process, which is the only channel through which information feeds back into
the work \eqref{eq:work}.

\subsection{Reduction to the control variable}

\begin{lemma}[Data-processing reduction]\label{lem:dpi}
Let $A:=a_{[0,\tau]}$ denote the control trajectory generated by the protocol.
Then
\begin{equation}\label{eq:IleIA}
I \;=\; \mathcal I\!\bigl(X_{[0,\tau]};\,Y\bigr)\;\ge\;
\mathcal I\!\bigl(X_{[0,\tau]};\,A\bigr),
\end{equation}
and the work \eqref{eq:work}, hence $\Wext$, is a functional of $(X,A)$ only.
Consequently the thermodynamically \emph{usable} information is
$I_A:=\mathcal I(X_{[0,\tau]};A)$, and
\begin{equation}\label{eq:IAbound}
I_A \;\le\; H(A),
\end{equation}
where $H(A)$ is the Shannon entropy of the control trajectory.
\end{lemma}

\begin{proof}
The control trajectory $A$ is, by definition of the protocol, a
(measurable, possibly randomised) function of the record $Y$ alone:
$A=g(Y,\zeta)$ with $\zeta$ a randomisation variable independent of $X$. Hence
the Markov chain
\[
X \longrightarrow Y \longrightarrow A
\]
holds (given $Y$, $A$ is conditionally independent of $X$). The data-processing
inequality for mutual information gives
$\mathcal I(X;A)\le \mathcal I(X;Y)=I$, which is \eqref{eq:IleIA}.

That $\Wext$ depends on $(X,A)$ only is immediate from \eqref{eq:work}: the
integrand $V(X_t)$ and the integrator $\dd a_t$ involve only $X$ and $A$. The
measurement record $Y$ influences the work \emph{exclusively} through the
control $A$ it produces; this is the precise sense in which $I_A$ is the
usable part.

Finally, $\mathcal I(X;A)=H(A)-H(A\mid X)\le H(A)$ because conditional entropy
is nonnegative; this is \eqref{eq:IAbound}.
\end{proof}

\eqref{eq:IleIA}--\eqref{eq:IAbound} say that, although the raw measurement may
carry arbitrarily much information $I$, only the part transduced into the
control can be converted to work, and that part is capped by the entropy of the
binary switching process.

We may in fact sharpen the second law \eqref{eq:2ndlaw} by replacing $I$ with
the smaller $I_A$. Indeed, the closed-loop dynamics \eqref{eq:langevin} and the
work \eqref{eq:work} depend on the record $Y$ \emph{only} through the control
process $A=g(Y,\zeta)$: conditionally on $A$, the particle evolution and the
work are independent of the residual measurement detail in $Y$. Hence the
control process $A$ is itself a (generally randomised, history-dependent)
``feedback variable'' that drives \eqref{eq:langevin}, and the
Sagawa--Ueda generalised second law applied to this effective measurement
outcome $A$ yields
\begin{equation}\label{eq:2ndlaw2}
\Wext \;\le\; -\Delta F + \kB T\, \mathcal I(X_{[0,\tau]};A)
\;=\;-\Delta F + \kB T\, I_A \;\le\; -\Delta F + \kB T\, H(A).
\end{equation}
This is consistent with, and stronger than, \eqref{eq:2ndlaw}, since
$I_A\le I$ by \eqref{eq:IleIA}; the gap $I-I_A\ge0$ is exactly the part of the
acquired information that is discarded in forming the binary control and is
thermodynamically inert. (If one prefers to invoke \eqref{eq:2ndlaw} verbatim
with the raw record $Y$, all statements below remain valid \emph{a fortiori}
with $I$ in place of $I_A$, because $I_A\le I\le$ the relevant entropy bounds;
the $I_A$ form is the tighter one.)

\subsection{Bounding the entropy of the control trajectory}

We now bound $H(A)$ in the two protocol classes.

\begin{theorem}[Maximal entropy reduction, discrete time]\label{thm:Imax-discrete}
Suppose the controller makes $N$ measurements and may switch the potential at
most once after each measurement (one binary decision per measurement, the
standard flashing-ratchet protocol). Then for every feedback protocol
\begin{equation}\label{eq:Imax-discrete}
I_A \;\le\; H(A) \;\le\; N\log 2 \;=:\; I_{\max}^{(N)} .
\end{equation}
The bound is tight: $H(A)=N\log 2$ holds for any protocol whose binary
decisions $D_k\in\{0,1\}$ (``set $a=$ ON or OFF on the $k$-th interval'') are
\emph{jointly independent} and uniform, and the \emph{usable} bound
$I_A=N\log2$ is attained, in the noiseless-measurement limit, by a deterministic
protocol whose decisions are simultaneously jointly independent and uniform
(for instance $D_k=\mathbf 1\{X_{v_k}-X_{u_k}>0\}$ read off disjoint
sub-windows; see the proof). It is \emph{not} attained by a position-threshold
``sign rule'' $D_k=\phi(X_{t_k})$, because the sampled states of the closed-loop
diffusion are correlated, so the resulting bits are correlated and
$H(A)<N\log2$ strictly.
\end{theorem}

\begin{proof}
The control trajectory is determined by the sequence of binary decisions
$D=(D_1,\dots,D_N)\in\{0,1\}^N$ (the value of $a$ on each inter-measurement
interval), because $a_t$ is piecewise constant and equals $D_k$ on
$[t_k,t_{k+1})$. Hence $A$ is a deterministic bijective recoding of $D$, so
$H(A)=H(D)$. By subadditivity of Shannon entropy and the fact that each $D_k$
is a binary random variable,
\[
H(D)\le\sum_{k=1}^N H(D_k)\le \sum_{k=1}^N \log 2 = N\log 2,
\]
using $H(D_k)\le\log 2$ with equality iff $D_k$ is uniform. With
\eqref{eq:IAbound} this proves \eqref{eq:Imax-discrete}.

\emph{Tightness of $H(A)=N\log2$:} if the $D_k$ are i.i.d.\ uniform on
$\{0,1\}$ then $H(D)=\sum_k H(D_k)=N\log2$, so the chain of inequalities is an
equality.

\emph{Achievability of $I_A=N\log 2$.} To saturate \eqref{eq:Imax-discrete}
both inequalities must be equalities: $H(A\mid X)=0$ \emph{and} $H(D)=N\log2$.
The first requires that $A$ (equivalently $D$) be a deterministic function of
the path $X$; the second, by the equality case of subadditivity, requires that
$D_1,\dots,D_N$ be \emph{jointly independent} and each uniform on $\{0,1\}$.
Both conditions can be met simultaneously by a causal, noiseless protocol,
because the diffusion path carries an inexhaustible amount of fine-scale
randomness that the controller can read off in disjoint time-windows. Concretely,
pick $N$ disjoint sub-intervals $[u_k,v_k]\subset[t_{k-1},t_k]$ (one strictly
before each decision instant) and set
\[
D_k \;=\; \mathbf 1\{X_{v_k}-X_{u_k}>0\}.
\]
Each $D_k$ is the sign of a Brownian-type increment of $X$ over a disjoint
window. By symmetry of the increment under the noiseless dynamics each $D_k$ is
balanced, $\Pr[D_k=1]=\tfrac12$; and increments over \emph{disjoint} time
windows of \eqref{eq:langevin} are independent (independence of Wiener
increments, with the drift over a vanishingly short window contributing
nothing in the limit), so the $D_k$ are jointly independent. Hence
$H(D)=\sum_k H(D_k)=N\log2$, while $D=A$ is a deterministic functional of $X$,
giving $H(A\mid X)=0$ and therefore
\[
I_A=\mathcal I(X;A)=H(A)-H(A\mid X)=H(D)=N\log2 .
\]
This proves $I_{\max}^{(N)}=N\log2$ is attained, so the bound
\eqref{eq:Imax-discrete} is tight.

\emph{Caveat on the naive ``sign rule''.} It is \emph{not} true that the
familiar position-thresholding protocol $D_k=\phi(X_{t_k})$ (e.g.\ a sign rule
on the local force, sampled at the decision instants) saturates the bound, even
when each $D_k$ is made marginally balanced. The reason is that the sampled
states $X_{t_1},\dots,X_{t_N}$ of the closed-loop Markov diffusion are
\emph{correlated}, so the induced bits $D_k=\phi(X_{t_k})$ are correlated and
\[
H(D)\;=\;H\bigl(\phi(X_{t_1}),\dots,\phi(X_{t_N})\bigr)\;<\;\sum_k H(D_k)\;=\;N\log2
\]
strictly whenever consecutive samples are genuinely dependent. (A numerical
check on a balanced sign rule applied to an AR(1)/Ornstein--Uhlenbeck sampled
chain gives, e.g., $H(D)\approx 2.19$ nats for $N=4$ against $N\log2\approx2.77$
nats.) Thus the per-decision cap of $\log2$ nats is a genuine upper bound,
attained only by protocols whose successive decisions are \emph{informationally
fresh}, i.e.\ jointly independent; reading correlated samples leaves part of the
$N\log2$ budget unused.

\emph{Effect of measurement noise.} If the measurement is noisy,
$D_k=\phi(Y_k)$ with $Y_k=m(X_{t_k},\xi_k)$ and $\xi_k$ independent of $X$, then
$D$ is no longer $\sigma(X)$-measurable, $H(A\mid X)>0$, and by the strict
data-processing gap $I_A<H(A)\le N\log2$: noise lowers the usable information
below the cap.

\emph{Information optimum $\ne$ work optimum.} We stress that the
disjoint-increment construction above maximises the \emph{information} $I_A$
but does \emph{not} by itself extract any work. Indeed, if the windows
$[u_k,v_k]$ are read during OFF (free-diffusion) phases, the increments
$X_{v_k}-X_{u_k}$ are statistically independent of the particle's phase
relative to the potential, so the resulting switching is uncorrelated with
$V(X)$ at the switch instants and the mean work \eqref{eq:work} vanishes,
$\Wext=0$, even though $I_A=N\log2$. Thus this protocol shows the
\emph{information} bound $I_{\max}=N\log2$ is attained, while leaving the
\emph{work} bound $\Wext\le\kB T I_A$ extremely slack. Saturating the work
bound at positive work is a strictly stronger requirement, treated separately
in \S4.4: it demands that the $N$ independent balanced bits each arise from a
\emph{reversible, work-extracting} sub-cycle, not merely from any source of
fresh path randomness.
\end{proof}

\begin{theorem}[Maximal entropy-reduction \emph{rate}, continuous time]
\label{thm:Imax-cont}
Model a continuous-time controller with a finite actuation resolution: there is
a minimal dwell time $\delta>0$ and the controller may only \emph{begin or end}
an ON phase at the times of the grid
$\mathcal G_\delta=\{0,\delta,2\delta,\dots\}\cap[0,\tau]$ (any physical
actuator has finite switching speed and finite timing resolution; we make both
explicit). Set $m:=\lfloor\tau/\delta\rfloor$. Under this model the control
trajectory $A$ takes \emph{discretely many} values---it is constant on each of
the $m$ slots $[\,j\delta,(j+1)\delta)$ and equal to a bit $c_j\in\{0,1\}$
there---so its Shannon entropy is well defined and
\begin{equation}\label{eq:Icont2}
I_A\;\le\;H(A)\;=\;H(c_0,\dots,c_{m-1})\;\le\;\sum_{j=0}^{m-1}H(c_j)\;\le\; m\,\log 2
\;=\;\Bigl\lfloor\tfrac{\tau}{\delta}\Bigr\rfloor\log 2 .
\end{equation}
In particular the entropy-reduction \emph{rate} is finite,
\begin{equation}\label{eq:Icont}
\dot I_{\max}\;:=\;\sup_{\text{protocols}}\frac{I_A}{\tau}\;\le\;\frac{\log 2}{\delta},
\end{equation}
i.e.\ the controller can reduce entropy by at most $\log2$ per dwell time
$\delta$ (equivalently $\kB\log2$ of entropy per actuation slot). If one further
budgets the \emph{expected number of switches} by $\E[K]\le R\tau$ with
$R<\infty$ (each switch costs actuation work, so finite power forbids unbounded
switching), then only at most $R\tau$ of the $m$ slot-bits can differ from their
predecessor, and a sharper counting bound holds:
\begin{equation}\label{eq:Icont3}
I_A\;\le\;H(A)\;\le\;\log2 + H(K) + \E[K]\,\log\frac{e\,m}{\E[K]}
\;\le\;\log2 + H(K) + R\tau\,\log\frac{e\,m}{R\tau},
\end{equation}
where the last form uses $\E[K]\le R\tau\le m$ and monotonicity of
$x\mapsto x\log(em/x)$ on $(0,m]$. In the idealised zero-dwell, infinite-rate
limit $\delta\to0$ (no finite actuation budget) the bound \eqref{eq:Icont2}
diverges, reflecting that an idealised continuous noiseless controller with
unlimited actuation could in principle reduce entropy without bound; \emph{any}
finite resolution $\delta$ or switching rate $R$ renders the entropy reduction
and its rate finite.
\end{theorem}

\begin{proof}
\emph{The basic slot bound \eqref{eq:Icont2}.} By the resolution model the
control is constant on each slot $S_j:=[\,j\delta,(j+1)\delta)$, $0\le j\le m-1$
(the final partial slot, if any, inherits the value of slot $m-1$ and adds no
new degree of freedom). Writing $c_j\in\{0,1\}$ for the value of $a$ on $S_j$,
the trajectory $A$ is the deterministic bijective recoding of the finite bit
string $(c_0,\dots,c_{m-1})\in\{0,1\}^m$; hence $A$ is a discrete random
variable and $H(A)=H(c_0,\dots,c_{m-1})$ is an ordinary Shannon entropy.
Because $A$ is discrete, the data-processing identity
$I_A=\mathcal I(X;A)=H(A)-H(A\mid X)\le H(A)$ used in
Lemma~\ref{lem:dpi} is valid (it needs $H(A)$ to be a Shannon, not differential,
entropy). Subadditivity of Shannon entropy gives
$H(c_0,\dots,c_{m-1})\le\sum_{j}H(c_j)\le m\log2$ since each $c_j$ is binary.
This proves \eqref{eq:Icont2}; dividing by $\tau$ and using
$m=\lfloor\tau/\delta\rfloor\le\tau/\delta$ gives \eqref{eq:Icont}.

\emph{The switch-budgeted bound \eqref{eq:Icont3}.} A right-continuous
piecewise-constant trajectory on the grid is determined by (i) its initial value
$a_0=c_0\in\{0,1\}$, (ii) the number of switches
$K=\#\{1\le j\le m-1: c_j\ne c_{j-1}\}$, and (iii) the subset of slot-boundaries
at which a switch occurs, a $K$-element subset of the $m-1$ interior boundaries;
given (i)--(iii) all $c_j$ are fixed by alternation. Hence the map
$(a_0,K,\{\text{switch slots}\})\mapsto A$ is injective and, by the chain rule
and subadditivity of Shannon entropy,
\[
H(A)\le H(a_0)+H(K)+H(\{\text{switch slots}\}\mid K),\qquad H(a_0)\le\log2.
\]
Given $K$, the switch-slot subset is one of at most $\binom{m-1}{K}\le\binom{m}{K}$
choices, so by the uniform (maximum-entropy) bound and Jensen's inequality
applied to the concave function $K\mapsto\log\binom{m}{K}$,
\[
H(\{\text{switch slots}\}\mid K)\le \E\!\Bigl[\log\binom{m}{K}\Bigr]
\le \E[K]\,\log\frac{e\,m}{\E[K]},
\]
using the standard estimate $\binom{m}{k}\le (em/k)^{k}$, i.e.\
$\log\binom{m}{k}\le k\log(em/k)$, together with monotonicity of
$k\mapsto k\log(em/k)$ on $(0,m]$ and Jensen. Imposing $\E[K]\le R\tau$ and the
monotonicity of $x\mapsto x\log(em/x)$ yields the final form of
\eqref{eq:Icont3}. The data-processing inequality $I_A\le H(A)$ holds throughout
because $A$ is discrete.

\emph{The idealised limit.} As $\delta\to0$ we have $m=\lfloor\tau/\delta\rfloor
\to\infty$, so the right side of \eqref{eq:Icont2} diverges; with no switch
budget the entropy reduction is unbounded, while any fixed $\delta>0$ (or any
fixed $R<\infty$ in \eqref{eq:Icont3}) keeps it finite.
\end{proof}

\begin{remark}[Interpretation of $I_{\max}$]
Theorems \ref{thm:Imax-discrete}--\ref{thm:Imax-cont} give the
ratchet-specific maximal entropy reduction: \emph{per binary control decision
the feedback can reduce the system entropy by at most $\log 2$} (one nat-valued
bit; restoring units, $\kB\log2$ of entropy). In discrete time with $N$
decisions, $I_{\max}=N\log2$; in continuous time with finite switching rate
$R$ and dwell time $\delta$, the entropy reduction stays finite and only
diverges in the unphysical zero-dwell limit. The binary nature of ``flashing''
is what caps the per-action information at one bit.
\end{remark}

\section{(c) Complete thermodynamic bounds: work, power, efficiency}

We now assemble the consequences. Let $A$ be the control process, $I_A$ its
usable information \eqref{eq:IleIA}, and recall \eqref{eq:2ndlaw2}.

\subsection{Extracted work}

\begin{theorem}[Tight work bound]\label{thm:work}
For every feedback protocol on the flashing ratchet,
\begin{equation}\label{eq:workbound}
\Wext \;\le\; -\Delta F + \kB T\,I_A \;\le\; -\Delta F + \kB T\,H(A),
\end{equation}
and in the cyclic/steady operation of a ratchet, where the potential and the
particle distribution return statistically to their starting condition over a
period so that $\Delta F=0$, this reads
\begin{equation}\label{eq:workbound-cyclic}
\Wext \;\le\; \kB T\, I_A \;\le\; \kB T\, H(A).
\end{equation}
In discrete time with $N$ binary decisions, $\Wext\le N\,\kB T\log2$. Equality
in the first inequality of \eqref{eq:workbound} holds iff the feedback-controlled
forward process is statistically reversible (the saturation condition of the
generalised second law); equality in the second holds iff $H(A\mid X)=0$, i.e.\
the control is a deterministic function of the trajectory (noiseless,
non-randomised feedback).
\end{theorem}

\begin{proof}
\eqref{eq:workbound} is \eqref{eq:2ndlaw2}; \eqref{eq:workbound-cyclic} is its
specialisation to $\Delta F=0$. The numerical bound follows from
Theorem~\ref{thm:Imax-discrete}. The first equality condition is the
equality condition of the Sagawa--Ueda second law \eqref{eq:2ndlaw}, namely
vanishing of the total (system $+$ controller) entropy production, i.e.\
reversibility of the controlled dynamics. The second equality condition is the
equality condition of \eqref{eq:IAbound}: $I_A=H(A)-H(A\mid X)=H(A)$ iff
$H(A\mid X)=0$, i.e.\ $A$ is $\sigma(X)$-measurable.
\end{proof}

\subsection{Output power}

The output power is the mean extracted work per unit time. For cyclic
operation with period $\tau$ and $N$ measurements per period:

\begin{corollary}[Power bound]\label{cor:power}
\begin{equation}\label{eq:power}
\mathcal P \;:=\; \frac{\Wext}{\tau}\;\le\;\frac{\kB T\,I_A}{\tau}
\;\le\;\frac{\kB T\,N\log2}{\tau}.
\end{equation}
Writing $f=N/\tau$ for the measurement frequency, the power is bounded by
$\kB T\,(\log2)\,f$, i.e.\ by the information \emph{rate} times $\kB T$. In
continuous time, by Theorem~\ref{thm:Imax-cont},
$\mathcal P\le \kB T\,H(A)/\tau\le \kB T\,(\log2)/\delta$, which stays
finite for any finite dwell time $\delta>0$ (and is further reduced under a
finite switching-rate budget $R$)
\end{corollary}

\begin{proof}
Divide \eqref{eq:workbound-cyclic} by $\tau$ and insert the relevant
$I_{\max}$. The continuous-time statement divides $H(A)$ from
Theorem~\ref{thm:Imax-cont} by $\tau$.
\end{proof}

\subsection{Efficiency}

For an information-driven ratchet the natural thermodynamic efficiency
compares the work extracted to the information ``fuel'' supplied, both in
energy units. Define the \emph{information efficiency}
\begin{equation}\label{eq:eta-def}
\eta_{\mathrm{info}} \;:=\; \frac{\Wext + \Delta F}{\kB T\, I_A}\, .
\end{equation}

\begin{corollary}[Efficiency bound]\label{cor:eff}
For every feedback protocol with $I_A>0$,
\begin{equation}\label{eq:eta}
\eta_{\mathrm{info}} \;\le\; 1,
\end{equation}
with equality iff the controlled forward process is reversible (the saturation
condition of \eqref{eq:2ndlaw}). If, in addition, one charges the controller
the raw acquired information $I\ge I_A$ rather than the usable part $I_A$, the
realised efficiency
$\tilde\eta:=(\Wext+\Delta F)/(\kB T\,I)\le I_A/I\le1$ is further degraded by
the information-transduction loss $I_A/I$, which equals $1$ iff the control
loses no information about $X$ (lossless transduction, $X\to Y\to A$ with no
data-processing gap).
\end{corollary}

\begin{proof}
By \eqref{eq:workbound}, $\Wext+\Delta F\le \kB T I_A$; dividing by
$\kB T I_A>0$ gives \eqref{eq:eta}. Equality is the equality case of
\eqref{eq:2ndlaw}. The refined statement uses $\Wext+\Delta F\le\kB T I_A\le
\kB T I$ and divides by $\kB T I$; the factor $I_A/I\le1$ is exactly the
data-processing ratio from Lemma~\ref{lem:dpi}, equal to $1$ iff
$\mathcal I(X;A)=\mathcal I(X;Y)$, i.e.\ no information about $X$ is lost in
producing the control.
\end{proof}

\subsection{An explicit work-saturating construction}

Before stating the optimality criterion we exhibit a protocol that
\emph{simultaneously} attains $I_A=N\log2$ \emph{and} $\Wext=\kB T\,N\log2$,
so that the bound \eqref{eq:workbound-cyclic} is genuinely tight at positive
work (and not merely as an information identity). The construction is a chain
of $N$ independent reversible Szilard sub-cycles embedded in the flashing
ratchet.

\begin{proposition}[Reversible work-extracting chain]\label{prop:construct}
Assume the OFF potential is flat and the ON potential per period can be loaded
quasi-statically (the standard idealisation in which the ON potential is ramped
on and off slowly compared with the relaxation time, so each ramp is
reversible). Partition $[0,\tau]$ into $N$ identical sub-intervals
$J_k=[t_{k-1},t_k]$, $t_k-t_{k-1}=\tau/N$, each long enough that the particle
fully relaxes within it. In each $J_k$ run the following reversible Szilard
sub-cycle:
\begin{enumerate}
\item[\textup{(1)}] Start from the OFF (uniform on a period) equilibrium.
\item[\textup{(2)}] Measure on which half of the period the particle lies,
$D_k=\mathbf 1\{X_{t_{k-1}}\in \text{right half}\}\in\{0,1\}$ (a noiseless,
balanced binary measurement: by the OFF equilibrium symmetry
$\Pr[D_k=1]=\tfrac12$).
\item[\textup{(3)}] Conditioned on $D_k$, quasi-statically insert and then
relax a barrier/biasing ON potential that confines the particle to the
measured half and then lets it expand reversibly back to the full period,
extracting the isothermal free-energy difference $\kB T\log2$ of a one-bit
Szilard expansion.
\item[\textup{(4)}] Remove the potential (OFF), returning to the starting
equilibrium so the sub-cycle is closed, $\Delta F=0$ over $J_k$.
\end{enumerate}
Then:
\textup{(i)} $D_1,\dots,D_N$ are jointly independent and balanced, so
$H(D)=N\log2$ and $A$ (the resulting control trajectory) is a deterministic
function of the measurement record, whence $I_A=\mathcal I(X;A)=N\log2$;
\textup{(ii)} each sub-cycle is reversible and extracts mean work
$\langle W_{\mathrm{ext}}\rangle_k=\kB T\log2$, so
$\Wext=\sum_{k=1}^N\langle W_{\mathrm{ext}}\rangle_k=\kB T\,N\log2$ with
$\Delta F=0$; hence \eqref{eq:workbound-cyclic} holds with \emph{equality},
$\Wext=\kB T I_A=\kB T\,N\log2$.
\end{proposition}

\begin{proof}
\emph{(i) Independence and balance.} Step (4) returns the particle to the OFF
equilibrium at the end of each $J_k$, and (by the assumed full relaxation
within $J_k$) the law of $X_{t_k}$ is the OFF equilibrium \emph{independently}
of the entire history on $J_1,\dots,J_k$. Hence the initial state $X_{t_{k-1}}$
of sub-cycle $k$ is drawn from the OFF equilibrium independently of
$D_1,\dots,D_{k-1}$, and $D_k=\mathbf 1\{X_{t_{k-1}}\in\text{right half}\}$ is
therefore independent of $D_1,\dots,D_{k-1}$ and balanced by the reflection
symmetry of the uniform OFF equilibrium about the period's midpoint. By
induction $D_1,\dots,D_N$ are jointly independent and each uniform, so
$H(D)=N\log2$. The control trajectory $A$ is the deterministic alternation of
the conditional Szilard ramps dictated by $(D_k)$, hence a deterministic
function of the (noiseless) record; thus $H(A\mid X)=0$ and
$I_A=\mathcal I(X;A)=H(A)=H(D)=N\log2$.

\emph{(ii) Reversibility and work.} Within $J_k$ the protocol is a textbook
one-bit Szilard engine operated quasi-statically: conditioned on the measured
half, the system performs an isothermal reversible expansion from a
half-period to the full period, whose extracted work equals the free-energy
drop $\kB T\log\!\frac{\text{(full period)}}{\text{(half period)}}=\kB T\log2$.
Quasi-static operation makes the sub-cycle reversible (zero entropy
production in system $+$ bath beyond the unavoidable feedback term), so the
generalised second law \eqref{eq:2ndlaw2} holds with equality at the
sub-cycle level: $\langle W_{\mathrm{ext}}\rangle_k=\kB T\,\mathcal
I(X_{J_k};D_k)=\kB T\log2$ with $\Delta F_k=0$. Summing the independent,
statistically identical sub-cycles gives $\Wext=\kB T\,N\log2$ and
$\Delta F=0$, i.e.\ equality in \eqref{eq:workbound-cyclic}.
\end{proof}

\begin{remark}
The crucial difference from the disjoint-increment construction of
Theorem~\ref{thm:Imax-discrete} is that there the fresh bits were read from
\emph{work-irrelevant} free-diffusion increments, giving $I_A=N\log2$ but
$\Wext=0$; here each fresh bit is the outcome of a \emph{reversible
work-extracting} measurement of the work-relevant coordinate (which half of the
period the particle occupies), so the same $N\log2$ of information is fully
transduced into $\kB T\,N\log2$ of work. Decorrelation of successive bits is
supplied by the relaxation/reset in step (4), which is exactly what a plain
position-threshold sign rule on an \emph{un-reset}, correlated Markov chain
lacks.
\end{remark}

\subsection{Saturating protocols}

\begin{proposition}[Identification of optimal protocols]\label{prop:opt}
A discrete-time feedback flashing ratchet saturates the work bound
\eqref{eq:workbound} at the maximal value $\Wext=-\Delta F+\kB T\,N\log2$
(equivalently saturates the power bound \eqref{eq:power} at
$\kB T(\log2)N/\tau$ and the efficiency bound \eqref{eq:eta} at
$\eta_{\mathrm{info}}=1$ with $I_A=N\log2$) if and only if it satisfies all
three of:
\begin{enumerate}
\item[\textup{(i)}] \emph{Maximal usable information}: the binary decisions are
deterministic in the (noiseless) measurement and \emph{jointly independent} and
balanced, $\Pr[D_k=1]=\tfrac12$ and $H(D)=N\log2$, giving $I_A=N\log2$. \emph{A
plain position-threshold sign rule $D_k=\phi(X_{t_k})$ on an un-reset chain
does not qualify}, since correlations between successive samples make
$H(D)<N\log2$;
\item[\textup{(ii)}] \emph{Work-relevant, losslessly transduced information}:
the chain $X\to Y\to A$ has no data-processing gap
($\mathcal I(X;A)=\mathcal I(X;Y)$) \emph{and} each bit informs a control
action correlated with $V(X)$ at the switch instants, so that the information is
about the work-relevant coordinate rather than inert path detail;
\item[\textup{(iii)}] \emph{Thermodynamic reversibility}: the controlled
forward process is statistically reversible, so the generalised second law
\eqref{eq:2ndlaw} holds with equality.
\end{enumerate}
The construction of Proposition~\ref{prop:construct} satisfies
\textup{(i)--(iii)} and hence realises all three saturations simultaneously;
conversely each of \textup{(i)--(iii)} is necessary.
\end{proposition}

\begin{proof}
\emph{Sufficiency.} If (i)--(iii) hold then $I_A=N\log2$ (by (i) and
Lemma~\ref{lem:dpi}, since determinism gives $H(A\mid X)=0$ and joint
independence gives $H(A)=N\log2$), and reversibility (iii) makes the first
inequality of \eqref{eq:workbound} an equality, so
$\Wext=-\Delta F+\kB T I_A=-\Delta F+\kB T\,N\log2$. Dividing by $\tau$
saturates \eqref{eq:power} and, with $\Delta F=0$, \eqref{eq:eta} becomes
$\eta_{\mathrm{info}}=1$. Proposition~\ref{prop:construct} exhibits a protocol
meeting (i)--(iii), so the simultaneous optimum is non-vacuous.

\emph{Necessity.} Suppose all three saturations hold. Saturating
$\Wext=-\Delta F+\kB T\,N\log2$ forces $I_A=N\log2$ (else
$\Wext\le-\Delta F+\kB T I_A<-\Delta F+\kB T N\log2$); by the equality case of
\eqref{eq:Imax-discrete}, $I_A=N\log2$ requires $H(A\mid X)=0$ and
$H(D)=N\log2$, i.e.\ deterministic, jointly-independent, balanced decisions,
which is (i). Saturating the first inequality of \eqref{eq:workbound} is by
Theorem~\ref{thm:work} exactly reversibility (iii). Finally, achieving positive
maximal work with $I_A=N\log2$ requires that the information be transduced
losslessly into the control \emph{and} that the control be correlated with
$V(X)$ at the switching instants (otherwise, as the disjoint-increment example
shows, one may have $I_A=N\log2$ yet $\Wext=0<\kB T N\log2$); this is (ii).
Hence (i)--(iii) are necessary.
\end{proof}

\begin{remark}[Physical reading]
Conditions (i)--(iii) make the optimal feedback flashing ratchet completely
explicit: \emph{measure as informatively as the binary action allows (one
balanced \emph{and informationally fresh}, i.e.\ jointly independent, bit per
decision), about the work-relevant coordinate, feed all of it back without
loss, and switch quasi-statically/reversibly}. Proposition~\ref{prop:construct}
realises this as a chain of $N$ independent reversible Szilard sub-cycles, each
extracting exactly $\kB T\log2$ of work per balanced bit; the power is then
$\kB T\log2$ times the measurement rate, and the information efficiency is
unity. The binary ``flashing'' constraint caps the per-action benefit at one
bit ($\kB T\log2$ of free energy / work), the Landauer/Szilard value; the
maximal entropy reduction of the device is therefore $I_{\max}=N\log2$
(discrete) and finite for any physical actuation budget (continuous).
\end{remark}


\section*{Summary}
\begin{itemize}
\item[(a)] The information of a feedback flashing-ratchet protocol is the
mutual information $I=\mathcal I(X_{[0,\tau]};Y)$ between trajectory and
record, Definition~\ref{def:I}, with the causal/non-Markovian decomposition
\eqref{eq:Ichain} (a transfer-entropy sum; a transfer-entropy rate in
continuous time). The thermodynamically usable part is
$I_A=\mathcal I(X_{[0,\tau]};A)\le I$ (Lemma~\ref{lem:dpi}).
\item[(b)] The maximal entropy reduction is set by the binary control:
$I_{\max}=N\log2$ in discrete time (Theorem~\ref{thm:Imax-discrete}), and
remains finite for any finite switching rate and dwell time in continuous time
(Theorem~\ref{thm:Imax-cont}); both are saturated by deterministic noiseless feedback whose successive
binary decisions are jointly independent and balanced (a plain
position-threshold sign rule does \emph{not} saturate, since correlations among
samples make $H(A)<N\log2$).
\item[(c)] Hence the tight bounds
$\Wext\le -\Delta F+\kB T I_A\le -\Delta F+\kB T N\log2$
(Theorem~\ref{thm:work}),
$\mathcal P\le \kB T I_A/\tau\le \kB T(\log2)N/\tau$
(Corollary~\ref{cor:power}), and $\eta_{\mathrm{info}}\le1$
(Corollary~\ref{cor:eff}). These are \emph{simultaneously} saturated at
positive work by the explicit chain of $N$ reversible Szilard sub-cycles of
Proposition~\ref{prop:construct}, which attains $\Wext=\kB T\,N\log2$ with
$\Delta F=0$; the abstract optimality criterion is
Proposition~\ref{prop:opt}. We emphasise that maximising the information
$I_A=N\log2$ is necessary but \emph{not} sufficient for saturating the work
bound at positive work: the maximal information must be carried by bits that are
both work-relevant and losslessly transduced, as the disjoint-increment
protocol (with $I_A=N\log2$ yet $\Wext=0$) demonstrates.
\end{itemize}

\end{document}
